\subsection{Extensions of the character ring of the symmetric group}\label{sec:near-s3}
Consider the rings with basis $1,g,r,x$ and products
\[
 g^2=1,\quad gr=r,\quad r^2=1+g+r,\quad gx=x,\quad rx=2x,
 \quad x^2=1+g+2r+\kappa x
\]
for $\kappa=6n$, $n\geq0$. The cases $n=0,1,2$ are respectively
$F_{2,3}$, $F_{6,59}$ and $F_{12,260}$.
None admits a pivotal categorification. Put $d=3n+\sqrt{9n^2+6}$;
the dimension vector is $(1,1,2,d)$. The formal codegrees are
$D,D',3,2$, and the respective dimensions of the four summands of $I(1)$ are
\[
 1,\quad 1+nd,\quad 4+2nd,\quad 6+3nd.
\]
The corresponding low restriction matrix is forced to be
\[
 L=\begin{pmatrix}1&1&1&1\\0&0&1&1\\0&0&1&2\\0&n&2n&3n\end{pmatrix}.
\]
Here is an explicit justification despite the rational dependence of $1,1,2$.
The first column is the central unit; all other first entries equal one.
The irrational coefficients force the last row. The first row of
$LL^{\mathsf T}$, equal to that of the full Gram matrix, is
$(4,2,3,6n)$. The dimension-$1+nd$ column consequently has middle entries
zero. If $a,b$ are the middle entries of the dimension-$4+2nd$ column, then
$a+2b=3$ and $0\leq a\leq2$, forcing $(a,b)=(1,1)$; the remaining
column has middle entries $(1,2)$.

The low mass in $I(g)$ is $10+5nd$, while
$D=12+6nd$. Thus $2\FPdim(I_{\rm low}(g))-D=8+4nd>0$, contradicting
the first-moment bound. The case $n=0$ is also excluded without a pivotal
assumption: its Frobenius--Perron dimension is the integer $12$, so every
categorification is pseudo-unitary \cite[Proposition~8.24]{ENO}.
For every $n$, the characters with codegrees $2$ and $3$ vanish on $x$.
The remaining two characters are the Frobenius--Perron character and its
Galois conjugate. Since the ring is self-dual, \cite[Lemma~\ref{B-lem:galois}]{Finite}
therefore removes the pseudo-unitary hypothesis from the pivotal exclusion.
